\documentclass[11pt]{article}
\usepackage[margin=1in]{geometry}
\usepackage{amsmath,amssymb,amsthm}
\usepackage[colorlinks=true,urlcolor=blue,linkcolor=blue,citecolor=blue]{hyperref}
\newtheorem{theorem}{Theorem}[section]
\newtheorem{lemma}[theorem]{Lemma}
\newcommand{\T}{\mathbb T}
\newcommand{\Z}{\mathbb Z}
\newcommand{\Q}{\mathbb Q}
\newcommand{\C}{\mathbb C}
\newcommand{\dd}{\mathrm d}
\title{Which observation times determine a circle jump measure?\\
\large A single-time inverse theorem for Dirichlet Li coefficients}
\author{}
\date{27 November 2026}
\begin{document}
\maketitle
\begin{abstract}
For a circle jump measure with integral multiplicities and a canonically
convergent drift, a single time slice of the convolution semigroup
determines the measure for every such input precisely when $(\pi/t)^2$ is
irrational. At any time two indistinguishable measures can differ only
on a finite cyclic subgroup. For primitive Dirichlet $L$-functions, every
positive time is determining. Using Dirichlet GRH, we prove that the
complex Fourier multipliers $e^{-t\lambda(n)}$ distinguish all finite
products of primitive factors with multiplicity. Distinct products give a
fixed discrepancy in ratios of three consecutive multipliers in every
sufficiently long block. The proof uses almost periodic Fourier series,
quadratic Gauss sums and a two-step shift of completed $L$-functions.
\end{abstract}

\section{The observation and its arithmetic application}

Write $\T=\{z\in\C:|z|=1\}$. An admissible multiset $E$ in
$\T\setminus\{1\}$ has positive integral multiplicities $m_E(z)$ and satisfies
\begin{equation}
 \sum_{z\in E}m_E(z)|1-z|^2<\infty,\qquad
 \lim_{\epsilon\downarrow0}
 \sum_{\substack{z\in E\\|1-z|\ge\epsilon}}m_E(z)\operatorname{Im}z
 \quad\text{exists}. \label{ad}
\end{equation}
Sums over $z\in E$ range over distinct points. The first condition makes
each cutoff finite. Define, for $n\in\Z$,
\begin{equation}
 \lambda_E(n)=\lim_{\epsilon\downarrow0}
 \sum_{\substack{z\in E\\|1-z|\ge\epsilon}}m_E(z)(1-z^n),
 \qquad H_{E,t}(n)=e^{-t\lambda_E(n)}. \label{obs}
\end{equation}
The limit exists because $1-z^n+in\operatorname{Im} z=O_n(|1-z|^2)$
near $1$. We have $H_{E,t}(0)=1$ and $H_{E,t}(-n)=\overline{H_{E,t}(n)}$.
Each cutoff defines a compound Poisson probability law on $\T$ with these
truncated Fourier multipliers, in the convention
$\widehat\mu(n)=\int z^n\,\dd\mu(z)$. Compactness of probability measures
on $\T$ and uniqueness of Fourier coefficients give a limiting law with
multipliers \eqref{obs}. The limit fixes the drift from the atoms. A time
slice means this entire law, or equivalently its frequency-labelled
coefficients at a single time.

\begin{theorem}[Determining times]\label{main}
Fix $t>0$. If $H_{E,t}=H_{F,t}$, then $E$ and $F$ agree with
multiplicity outside the roots of unity of some common finite order.
The map $E\mapsto H_{E,t}$ on admissible multisets is injective if and
only if $(\pi/t)^2\notin\Q$. At every exceptional time an
indistinguishable pair can be obtained by finite additions to any common
admissible multiset.
\end{theorem}

The zero-location theorem of the pinned Lean development \cite{vendor}
places each nontrivial zero $\rho$ of a primitive Dirichlet $L$-function
on $\operatorname{Re}s=1/2$. Its Cayley point $z=1-1/\rho$ therefore belongs
to $\T\setminus\{1\}$, with $z=-1$ for a central zero. The cutoff
$|1-z|\ge\epsilon$ means precisely $|\rho|\le1/\epsilon$.
The classical zero count gives $\sum_\rho|\rho|^{-2}<\infty$.
Starred convergence of the complex Li zero sum gives the second condition
of \eqref{ad} and identifies \eqref{obs} with the complex Li coefficients
in that convention \cite{lagarias}. For a finite product $P$ of primitive
Dirichlet $L$-functions, including the empty product, take the multiset
union $E_P$ and write $H_{P,t}=H_{E_P,t}$. The primitive conductor-one
factor is $\zeta$. The circle law in the merged article \cite{article}
uses $\operatorname{Re}\lambda(n)$ and thus symmetrizes the Cayley
points. The complex coefficients here retain the signs of the ordinates.

\begin{theorem}[Rigidity for primitive products]\label{rigid}
For any $t>0$, if $H_{P,t}(n)/H_{Q,t}(n)$ tends to a finite nonzero
complex limit as $n\to+\infty$, then $P,Q$ have the same primitive
factors with multiplicity, and the limit equals $1$. More precisely, put
\[
 C_{P,t}(n)=\frac{H_{P,t}(n+1)H_{P,t}(n-1)}{H_{P,t}(n)^2}.
\]
For distinct factor multisets there are $\eta>0$ and an integer $L\ge1$
such that every block of $L$ consecutive positive integers contains an
$n$ with $|C_{P,t}(n)/C_{Q,t}(n)-1|>\eta$.
\end{theorem}

At $t=\pi$, a single atom at $i$ and one at $-i$ give the same circle law
since $i^n-(-i)^n\in2i\Z$ for all $n$. At the same time every primitive
Dirichlet product is determined among all such products of arbitrary
degree and conductor. The discrepancy and the block length may depend
on both products and on $t$.

\paragraph{What the symmetric law determines.}
Write $\overline P$ for the product with every character conjugated.
Its Cayley multiset is the complex conjugate of $E_P$, so
$|H_{P,t}|^2=H_{P\overline P,t}$. Theorem~\ref{rigid} shows that even a
finite nonzero limiting ratio $|H_{P,t}(n)|/|H_{Q,t}(n)|$ determines
equality of the total multiplicities in each character conjugacy orbit.
Conversely those totals determine the symmetric law. For a product the
ambiguity is precisely independent conjugation of individual factors,
allowing permutation. For a single primitive character the original
circle kernel in \cite{article}, at any one time $t>0$, therefore
determines the character up to complex conjugation. Retaining the
complex phases of the positive-frequency multipliers distinguishes the
two orientations whenever the character is nonreal, by Theorem~\ref{rigid}.

\section{From infinitely many angles to finite torsion}

Set $\Delta^2h(n)=h(n+1)-2h(n)+h(n-1)$ and
$D(n)=\lambda_E(n)-\lambda_F(n)$. Direct calculation at finite cutoffs
gives, on passing to the limit,
\begin{equation}
 u(n):=\Delta^2D(n)=\int_{\T}z^n\,\dd\tau(z),\qquad
 \tau=\sum_z |1-z|^2\big(m_E(z)-m_F(z)\big)\delta_z . \label{diff}
\end{equation}
The measure has finite total variation. Thus $u$ is a uniform limit of finite
trigonometric sums on $\Z$.

\begin{lemma}\label{per}
For each $\delta>0$ the shifts $h\in\Z$ with
$\sup_n|u(n+h)-u(n)|<\delta$ have bounded gaps. If $e^{-tu(n)}=1$
for all $n$, $\tau$ is supported on a finite cyclic subgroup of $\T$.
If $e^{-tu(n)}$ differs from $1$ somewhere, its distance from $1$ exceeds
a positive constant somewhere in every sufficiently long block.
\end{lemma}
\begin{proof}
Approximate $u$ within $\delta/4$ by finitely many terms $\sum_j b_jz_j^n$.
In the compact closure of the orbit $(z_1^h,\ldots,z_m^h)$ take an identity
neighborhood on which $\sum_j|b_j|\,|z_j^h-1|<\delta/2$.
Finitely many translates by orbit elements cover the closure. Subtracting
their finitely many indices from any integer $k$ gives at least one
return, at bounded distance from $k$. Every return is a $\delta$-almost period.

When $e^{-tu}=1$, the values of $u$ belong to $\beta\Z$ for
$\beta=2\pi i/t$. A positive $|\beta|/2$-almost period $q$ is an exact
period. All Fourier coefficients of $(z^q-1)\tau$ then vanish. By density
of Laurent polynomials this measure is zero, giving the asserted support.
For the last assertion, the exponential is uniformly continuous on the
bounded closure of the range of $u$. Sufficiently accurate almost periods
translate any fixed discrepancy to a relatively dense set on which a
positive discrepancy persists.
\end{proof}

\section{The exceptional times}

\begin{proof}[Proof of Theorem~\ref{main}]
When two observations agree, $D(n)\in\beta\Z$ and hence $u(n)\in\beta\Z$.
By Lemma~\ref{per} and $|1-z|^2>0$ for $z\ne1$, the multiplicities can
differ only at $q$th roots of unity for some $q$. The finite modifications
now give the entire difference without any residual linear term,
\[
 D(n)=\sum_{j=1}^{q-1}b_j(1-\zeta_q^{jn}),\qquad
 b_j\in\Z,\quad \zeta_q=e^{2\pi i/q}.
\]
If $D$ is zero, \eqref{diff} and Fourier uniqueness imply $E=F$.
Otherwise $D(n_0)=m\beta$ with $m\ne0$, placing $\beta$ in
$K=\Q(\zeta_q)$. Every automorphism $\sigma_a:\zeta_q\mapsto\zeta_q^a$
satisfies $\sigma_a D(n)=D(an)$, so $\sigma_a(\beta)/\beta$ is a
nonzero rational. These ratios form the image of a homomorphism from a
finite group to $\Q^\times$ since each ratio is rational. They all equal
$\pm1$, so $\beta^2\in\Q$. This proves necessity for a collision.

Conversely, if $\beta^2\in\Q$, the imaginary quadratic field
$\Q(\beta)$ has fundamental discriminant $d<0$ and
$\beta/\sqrt d\in\Q^\times$ since both elements have trace zero. Take
the primitive odd quadratic character $\chi_d$ modulo $q=|d|$ associated
to this field \cite{conrad}. Its additive Fourier transform is
\[
 G(n)=\sum_{j\bmod q}\chi_d(j)\zeta_q^{jn}
     =\chi_d(n)G(1),\qquad G(1)^2=-q=d .
\]
Indeed a change of variable proves the transform formula for unit $n$.
When $g=(n,q)>1$, multiplication by every element of the kernel of
reduction from units modulo $q$ to units modulo $q/g$ fixes the additive
character. Primitivity gives a kernel element on which $\chi_d$ is
nontrivial, forcing $G(n)=0$. Finite Parseval gives $|G(1)|^2=q$,
and oddness gives $\overline{G(1)}=-G(1)$.

Take a positive integer $c$ with $cG(1)/\beta\in\Z$. To one multiset
add the positive part of $c\chi_d(j)$ at each $\zeta_q^j$, and to the
other add its negative part in absolute value. Neither addition charges
$1$. Their difference gives $D(n)=-cG(n)\in\beta\Z$ since
$\sum_j\chi_d(j)=0$. These distinct multisets have identical
observations. Any common admissible background can be added.
\end{proof}

\section{Why primitive Euler products have no exceptional times}

\begin{lemma}\label{finite}
If the nontrivial zero multisets of two finite primitive Dirichlet
products $P,Q$ differ in finitely many zeros counted with multiplicity,
the primitive factor multisets agree.
\end{lemma}
\begin{proof}
For conductor $q$ and parity $a\in\{0,1\}$, complete $L(s,\chi)$ as
$(q/\pi)^{(s+a)/2}\Gamma((s+a)/2)L(s,\chi)$, multiplied by $s(s-1)$
for the conductor-one factor. Each is entire of order at most one with
zero divisor exactly the nontrivial zeros \cite{lagarias}.
Let $\Xi_P,\Xi_Q$ be their products. Hadamard factorization gives
\[
 \frac{\Xi_P(s)}{\Xi_Q(s)}=e^{a_0+b_0s} R(s)
\]
with $R$ a nonzero rational function, since the common genus-one
canonical factors cancel. Divide the relation at $s+2$ by its value at
$s$. Every gamma argument shifts by one, giving the rational ratio
$\Gamma(w+1)/\Gamma(w)=w$. All completion factors have rational shift
quotients up to nonzero constants. It follows that for $F=P/Q$, the
function $F(s+2)/F(s)$ is rational.

Along the positive real axis $L(\sigma,\chi)=1+O(2^{-\sigma})$,
by bounding $\sum_{n\ge2} n^{-\sigma}$. Thus $F(\sigma)$ and
$F(\sigma+2)/F(\sigma)$ equal $1+O(2^{-\sigma})$. A nonzero rational
function cannot decay exponentially at positive infinity, by its
leading Laurent term. Therefore $F(s+2)/F(s)=1$ identically. Iterating
the period along the real axis gives $F(\sigma)=1$ for every $\sigma>1$,
hence $P=Q$ by continuation.

Take a common multiple $M$ of all conductors. Uniqueness of Dirichlet
series at the prime coefficient $p\nmid M$ gives
$\sum_{\chi\in P}\chi(p)=\sum_{\psi\in Q}\psi(p)$.
Each reduced residue class modulo $M$ contains a prime, so these
character sums agree on $(\Z/M\Z)^\times$. Character orthogonality
recovers their multiplicities. Every induced character has a unique
primitive antecedent, which recovers all the given factors.
\end{proof}

\begin{proof}[Proof of Theorem~\ref{rigid}]
Take $E=E_P$ and $F=E_Q$. Then $C_{P,t}/C_{Q,t}=e^{-tu}$.
If this equals $1$ identically, Lemma~\ref{per} localizes the Cayley
multiplicity difference to finitely many points. The inverse
$\rho=1/(1-z)$ gives a finite difference of nontrivial zero divisors.
Lemma~\ref{finite} gives equality of all primitive factors. For
unequal products, Lemma~\ref{per} now yields the recurrent separation.
A nonzero limit of the multiplier ratio would force the ratio of the
three consecutive multipliers to tend to $1$, contrary to that
separation. Equal primitive products give ratio $1$.
\end{proof}

\paragraph{Relation to earlier work.}
Lagarias establishes starred convergence and asymptotics for automorphic
Li coefficients \cite{lagarias}. The merged article develops almost
periodic second differences to distinguish unexponentiated Li sequences
and constructs circle kernels from their real parts \cite{article}.
Nickl and S\"ohl obtain identifiability for a finite-activity circle
model under an intensity bound allowing a single logarithm branch
\cite{nickl}. Godsil and Zhang characterize square-root spectra of
finite periodic oriented and signed Cayley graphs \cite{godsil}.
Here equality of exponentiated data for possibly infinite integral
multisets forces the ambiguity onto one finite cyclic group. Quadratic
Gauss sums give all exceptional times, and the primitive Euler product
eliminates them. The scalar heat trace sums over frequencies and supplies
a different inverse question.

\begin{thebibliography}{9}\small
\bibitem{article} \emph{Beyond the Riemann hypothesis for Dirichlet
$L$-functions: consequences, transfers and limits}, section on Li
coefficients, moment pencils and positivity certificates.
\href{https://github.com/VladimirReshetnikov/ProveIt/blob/b6253012cf30c80a989c871906e097a27f23a27c/Analysis/GeneralizedRiemannHypothesis/Article/sections/li.tex}{Pinned article source}.
\bibitem{vendor} \emph{The generalized Riemann hypothesis for Dirichlet
$L$-functions, in Lean 4}, readable edition, upstream revision
\texttt{72fead102f66}.
\href{https://github.com/VladimirReshetnikov/ProveIt/blob/b6253012cf30c80a989c871906e097a27f23a27c/Analysis/GeneralizedRiemannHypothesis/vendor/grh-lean/grh/src/Main.lean}{Vendored theorem declarations}.
\bibitem{lagarias} J.~C.~Lagarias, Li coefficients for automorphic
$L$-functions, \emph{Ann. Inst. Fourier} \textbf{57} (2007), 1689--1740.
\href{https://doi.org/10.5802/aif.2311}{doi:10.5802/aif.2311}.
\bibitem{conrad} B.~Conrad, \emph{Quadratic characters associated to
quadratic fields}, Stanford Math 248A.
\href{https://math.stanford.edu/~conrad/248APage/handouts/quadchar.pdf}{Author's notes}.
\bibitem{nickl} R.~Nickl and J.~S\"ohl, Bernstein--von Mises theorems for
statistical inverse problems II: compound Poisson processes,
\emph{Electron. J. Statist.} \textbf{13} (2019), 3513--3571.
\href{https://doi.org/10.1214/19-EJS1609}{doi:10.1214/19-EJS1609}.
\bibitem{godsil} C.~Godsil and X.~Zhang, Oriented or signed Cayley graphs
with all eigenvalues integer multiples of $\sqrt\Delta$.
\href{https://arxiv.org/abs/2405.14140}{arXiv:2405.14140}.
\end{thebibliography}
\end{document}
